\section{Benchmark}
\label{sec:benchmark}


We construct a map-aware benchmark for learning the correspondence among an overhead tactical radar map, a first-person view (FPV), and the camera pose in Counter-Strike 2 (CS2).
Each example contains a map-specific radar image, an FPV image captured from the running game, and a five-dimensional (5-D) pose consisting of two horizontal coordinates, a vertical coordinate, pitch, and yaw.
The benchmark supports both directions of the correspondence: localizing the camera from an FPV/radar pair and synthesizing an FPV image from a radar map and a specified pose.
It also contains ordered short trajectories, allowing image generation to be evaluated on temporally continuous motion rather than only on independently sampled frames.

The construction is a semi-automatic pipeline with manual review: game capture produces paired visual and state observations; map resources provide the radar coordinate system; preprocessing reconciles the two state streams; quality control fixes the retained source samples and per-map height ranges; and the final split is formed from complete map records.
Figure~\ref{fig:cs2-benchmark-pipeline} summarizes this pipeline.

\begin{figure*}[t]
    \centering
    \fbox{\parbox[c][3.35cm][c]{0.94\linewidth}{\centering
    \textit{Placeholder for the benchmark construction figure.}\\[2pt]
    The final figure should show, from left to right, the controlled CS2
    setup and the three capture streams (FPV screenshot, GSI state, and
    process-memory state), radar extraction and map-coordinate conversion,
    Mem/GSI consistency filtering, missing-image removal, and audit,
    full-corpus height calibration, \texttt{file\_num}-record pool assignment,
    and the final localization and conditional-generation tasks.  It should
    distinguish individual-frame files from ordered 64-frame continuous
    clips.}}
    \caption{Placeholder.  The final figure should visualize the complete
    data path and identify the artifact produced at each stage, including
    the point at which record-level isolation is enforced.}
    \label{fig:cs2-benchmark-pipeline}
\end{figure*}


\subsection{Data Acquisition and Canonical Representation}

\paragraph{Controlled CS2 capture.}
The data are collected from a CS2 client in a controlled practice mode at a target rate of 16 frames per second.
We manually launch the game from Steam with \texttt{-console -insecure} and use a windowed 4:3, 1024$\times$768 display.
We use console commands to hide the selected HUD overlays, weapon and arm models, and crosshair so that the recorded FPV image exposes the game scene without those interface elements.
The recorder observes the selected player from the spectator view while the game is running.

Under the official Game State Integration (GSI) configuration, the CS2 client sends game state updates to a local endpoint, together with the relevant map and player state.
GSI provides the player position and game metadata but does not expose the complete camera viewpoint.
Therefore, the recorder reads the running game process to obtain the position and the pitch/yaw viewpoint fields from the CS2 process memory.
In each capture-loop iteration, it associates a whole-window screenshot with the most recent GSI payload and a process-memory readout. These observations are then paired by the recorder loop.
% the implementation does not provide a shared hardware timestamp or claim timestamp-atomic synchronization.

Raw observations are stored in \texttt{record\_N.npy} batches.
Each entry contains the FPV image path and the memory and GSI dictionaries, while the images are retained as individual JPEG files.
The resulting identity, \texttt{file\_numN\_frame\_M}, preserves both the record-batch identifier and the frame identifier.
Here, a record denotes one \texttt{record\_N.npy} batch identified by \texttt{file\_num}; it is not a physical capture-session identifier.
The complete map--\texttt{file\_num} group is carried through preprocessing and later serves as the atomic split unit.

\paragraph{Radar maps and coordinate transform.}
The radar image for each map is obtained from the CS2 resource package using Source 2 Viewer. The corresponding map-overview metadata supplies an offset and scale that align game-world positions with the radar image.
Let $(x_w,y_w,z_w)$ denote the camera position, where $z_w$ is the camera eye height, and let $(offset_x, offset_y, scale)$ denote the map-specific offset and scale.
We convert the position to the canonical radar coordinates
\begin{equation}
    x_r = \frac{x_w-offset_x}{scale}, \qquad
    y_r = \frac{offset_y-y_w}{scale}, \qquad
    z_r = \frac{z_w}{scale}.
\end{equation}
The reversal in the second equation accounts for image coordinates.
The implementation casts the transformed coordinates to integer radar units; $x_r$ and $y_r$ are expressed in the coordinate system of a 1024$\times$1024 radar image, while $z_r$ is retained as the vertical coordinate.
The upper- and lower-floor radar layers are blended for \texttt{de\_nuke}, \texttt{de\_train}, and \texttt{de\_vertigo}, as specified by the multi-layer map assets.

\paragraph{State reconciliation and sample representation.}
The preprocessing stage independently transforms the memory position and the GSI position into radar coordinates.
A frame is removed when the Euclidean distance between the two horizontal positions exceeds 3 radar pixels.
A frame whose referenced FPV image is missing is also removed.
For every retained frame, the image is assigned a stable name and the metadata are written to the map's \texttt{positions.json}.
A canonical sample is
\begin{equation}
    (x_r,y_r,z_r,\theta_h,\theta_v,\mathrm{map},\mathrm{file\_frame}),
\end{equation}
where $\theta_h$ is yaw and $\theta_v$ is pitch, both stored in radians.
Thus, a sample identifies one FPV image, one map-specific radar image, and one 5-D pose.


\subsection{Quality Control}
\label{subsec:data_filter}

\paragraph{Integrity review.}
The construction first checks frame identities for duplicates and validates finite pose values, map labels, source images, and radar assets.
% After that, we filter frames when the Euclidean distance between the GSI and process-memory positions exceeds a fixed threshold of 3 pixel units, accounting for local-machine freezes and network fluctuations.
It then reports coordinate candidates using out-of-window coordinates and large adjacent-frame changes in position or angle.
These conditions are manual review triggers rather than unconditional rejection rules, because multi-layer maps and camera motion can produce large valid changes.
We first identify large changes in horizontal and vertical position and viewing angle and compare them with changes that can arise from player movements within the game.
Large but valid changes in position and viewing angle can occur during movement on slopes, rapid turning, jumping, crouching, and falls from normal heights.
Afterwards, we manually review samples that exceed the large valid change ranges and reject unsuitable samples from our dataset.
Rejected large changes in horizontal and vertical position are mostly associated with random initial positions at the start of the game, respawning at a random location after death, and restarting after the time limit for a single game is exceeded.
Meanwhile, sudden changes in viewing angle appear to occur mostly when the player dies from a fall beyond the map boundaries or from a height greater than normal, while the recorded image switches from the first-person spectator view to a third-person view of the deceased body.

% The audit reports 10,345 coordinate candidates and no integrity-invalid samples.
% The manual review retains 9,705 candidate samples under the declared default action and excludes 18,756 samples, including 19 complete map records for conservative data quality.
% After the review, the source corpus contains 3,087,392 retained samples in 3,430 map records across the 14 maps.
% Of these samples, 2,280,253 belong to the ten seen maps and 807,139 to the four held-out maps.

\paragraph{Vertical normalization.}
After the review and before any record is assigned to a split, we compute an exact minimum and maximum of $z_r$ independently for every map over the retained source corpus. This provides a map-specific normalization of the vertical coordinate without relying on a global vertical map asset.
The task representation uses
\begin{equation}
    z' = \frac{z_r-z_{\min,m}}{z_{\max,m}-z_{\min,m}},
\end{equation}
where $m$ is the map. Because the calibration is computed before splitting, the same map-specific convention is used for training, few-shot learning, evaluation, and continuous samples.
% In particular, the supplied calibration for a held-out map is part of the benchmark input; held-out-map results should therefore be interpreted as zero-shot model transfer and few-shot adaptation under fixed map-level calibration, not as calibration-free transfer.

\paragraph{Record-level split construction.}
The complete map--\texttt{file\_num} group is the atomic split unit.
All frames from one such record are assigned to one pool before any frame sampling, preventing samples from the same record from entering different pools.
The record-pool ratios are 60\%/10\%/20\%/10\% for Seen-10 train/validation/discrete-test/continuous and 20\%/60\%/20\% for CrossMap-4 support/query/continuous.

For Seen-10 test selection, near-pose filtering removes candidates close to the selected training or validation poses; for CrossMap-4 query selection, it removes candidates close to the union of all five preconstructed support draws. The tolerances of near-pose filtering are 8 radar units in horizontal-plane Euclidean distance, 8 units in the vertical coordinate, and 5 degrees for each angle, with circular distance for yaw.
Eligible frames for every discrete split are then ordered within each record, thinned to a minimum gap of five frames, and selected by deterministic coverage sampling over $4\times4\times3\times8\times3$ quantile bins in $(x_r,y_r,z_r,\theta_h,\theta_v)$.

Continuous evaluation clips are selected only from the continuous record pool.
Each map contributes 20 clips, each clip contains exactly 64 ordered frames, and successive frame identifiers differ by one or two.
No frame is reused within the continuous evaluation split, and at most one clip is selected from a record.


\subsection{Dataset Splits and Statistics}
\label{subsec:benchmark-splits}

We divide the 14 map sources into a 10-map seen set and a 4-map held-out set according to differences in environment, lighting, texture, art style, and time of day.
In this setting, Seen-10 contains \texttt{cs\_agency}, \texttt{cs\_italy}, \texttt{de\_ancient}, \texttt{de\_anubis}, \texttt{de\_dust2}, \texttt{de\_inferno}, \texttt{de\_mirage}, \texttt{de\_nuke}, \texttt{de\_overpass}, and \texttt{de\_train}.
CrossMap-4 contains \texttt{cs\_office}, \texttt{de\_golden}, \texttt{de\_palacio}, and \texttt{de\_vertigo}.

For every Seen-10 map, the final discrete split contains 5,000 training frames, 500 validation frames, and 2,000 discrete test frames.
% Across the ten maps this gives 50,000 training frames, 5,000 validation frames, and
% 20,000 discrete test frames.
The continuous split contributes 20 clips of 64 frames per map, or 1,280 frames per map and 12,800 frames in total.

For every CrossMap-4 map, each preconstructed few-shot support draw contains 100 frames, the fixed discrete evaluation split contains 2,000 frames, and the continuous split contains 20 clips of 64 frames.
A single support draw therefore contains 400 frames over the four maps, while the discrete set contains 8,000 frames and the continuous set contains 5,120 frames.
The union of support records is disjoint from the discrete and continuous evaluation records, which are also disjoint from one another.
% Counting all five support draws separately gives 102,920 manifest row references; overlap among support draws yields 102,780 unique selected FPV images.

The pre-review source snapshot contains 3,106,148 preprocessed samples and 3,449 source records. After quality control, 3,087,392 (3M) samples and 3,430 records remain.
Specifically, Seen-10 contributes 2,280,253 retained samples and CrossMap-4 contributes 807,139.
The map-level source and retained counts are reported together with the map-specific $z$ extrema because maps differ in the number of records and vertical span.
Tables~\ref{tab:benchmark-statistics} and~\ref{tab:benchmark-splits} summarize the map-level corpus and the final task splits, respectively.

\begin{table*}[t]
    \centering
    \fbox{\parbox[c][4.05cm][c]{0.94\linewidth}{\centering
    \textit{Placeholder for the map-level dataset statistics table.}\\[2pt]
    The final table should contain one row for each of the 14 maps and report
    its Seen-10/CrossMap-4 role, preprocessed rows, approved-retained rows,
    source and retained \texttt{file\_num} records, frozen $z_{\min}$ and
    $z_{\max}$, and the exact radar asset (including blended assets).  An
    aggregate row should report 3,106,148 source rows, 3,087,392 retained
    rows, 3,449 source records, and 3,430 retained records.}}
    \caption{Placeholder.  The final table should make map imbalance,
    retained-corpus size, vertical calibration ranges, and radar-asset
    choices auditable.}
    \label{tab:benchmark-statistics}
\end{table*}

\begin{table*}[t]
    \centering
    \fbox{\parbox[c][3.75cm][c]{0.94\linewidth}{\centering
    \textit{Placeholder for the final split and task table.}\\[2pt]
    The final table should report, for Seen-10 and CrossMap-4, the map role,
    record-pool ratio, number of records, discrete rows or support/query
    rows, number of continuous clips, frames per clip, preconstructed support
    draws, the support draw used in the current experiments, and
    the exact task/evaluation split.  It should explicitly state the
    5,000/500/2,000 Seen-10 quotas, the 100-shot-per-map CrossMap quota, the
    five preconstructed support draws, the single \texttt{seed\_0} draw used in the
    current experiments, and the fixed 64-frame continuous clips.}}
    \caption{Placeholder.  The final table should connect each split to its
    intended task and document all quotas and disjointness constraints.}
    \label{tab:benchmark-splits}
\end{table*}


\subsection{Tasks and Evaluation Protocol}
\label{subsec:benchmark-tasks}

\paragraph{Task inputs and outputs.}
The benchmark defines two tasks: localization and generation, with two frame-level directions and one sequence-level evaluation setting.
For localization, the inputs are the FPV image, the corresponding overhead radar image, and the map identity; the output is the 5-D pose $(x,y,z,\mathrm{pitch},\mathrm{yaw})$.
For discrete conditional generation, the inputs are the radar image, map identity, and a specified pose; the output is the corresponding FPV image.
Continuous conditional generation uses the same conditioning variables for each frame but evaluates the generated images in the exact ordered 64-frame clips.

The pose variables are $(x_r,y_r,z_r,\theta_h,\theta_v)$, with yaw represented by $\theta_h$ and pitch represented by $\theta_v$.
The task serialization uses the normalized 5-D tuple
\begin{equation}
    \left(\frac{x_r}{1024},\frac{y_r}{1024},z',
    \frac{\theta_v}{2\pi},\frac{\theta_h}{2\pi}\right).
\end{equation}
The prompt convention defines yaw zero as east with clockwise increase, pitch zero as looking straight down, and pitch $\pi$ as looking straight up.

\paragraph{Seen-10 evaluation.}
The base model is trained only on the Seen-10 discrete train split.
Seen-10 validation records are reserved for checkpoint selection and hyperparameter decisions; they are not used to report the final test result.
Localization, discrete generation, and continuous generation are evaluated on their respective test splits.
The discrete test samples and continuous clips are record-disjoint from training and validation under the split construction above.

\paragraph{CrossMap-4 transfer.}
CrossMap-4 maps are excluded from base-model optimization.
Zero-shot transfer evaluates localization and discrete generation on the fixed discrete samples and evaluates continuous generation on the fixed continuous clips.
In the few-shot localization setting (100, 50, 20, 10), each shot-count run is independently initialized from the same Seen-10 checkpoint using nested subsets from the single \texttt{seed\_0} support draw.
Localization adaptation uses exactly 100, 50, 20, or 10 labeled frames per held-out map and is evaluated on the fixed CrossMap-4 discrete set, whereas generation adaptation uses the 100-shot \texttt{seed\_0} support draw and is evaluated on the corresponding discrete and continuous sets.
Each adapted model is also reevaluated on the corresponding Seen-10 split to assess retention.

All results are reported per map and as an equal-map macro average; few-shot results use a single support draw and do not include a support-selection confidence interval.
For sequence metrics, the 64-frame clip boundaries are preserved.
